% ============================================================
%  Abstract
% ============================================================
\begin{abstract}
Named entity recognition (NER) systems built on subword-tokenized language
models are brittle to the character-level noise ubiquitous in real-world
text---OCR errors, keyboard typos, transliteration artefacts, and deliberate
obfuscation.  We introduce \LightRet{}, a lightweight, vocabulary-free NER
model that achieves competitive accuracy on clean text while remaining robust
under adversarial character noise.

\LightRet{} is trained through a \emph{three-stage progressive distillation
pipeline}: (1)~sentence-level knowledge distillation from \textsc{BERT} into
a word-level \RetBERT{} student whose inputs are character embeddings;
(2)~token-level compression from \RetBERT{} into a compact
\textsc{BiGRU}--Transformer backbone; and (3)~noisy-student NER fine-tuning
with stochastic character-level augmentation and dynamic BIO label projection.

At its core, \LightRet{} uses \RetVec{}~\citep{sander2023retvec}---a frozen,
pretrained character-level embedder---eliminating vocabulary constraints
entirely.  On CoNLL-2003 English NER, \LightRet{} achieves
\textbf{\todo{XX.X}~F1} on clean text and \textbf{\todo{XX.X}~F1} under 10\%
character substitution noise, outperforming BERT-base by
\textbf{\todo{+X.X}~F1} in the noisy setting while being
\textbf{\(\approx28{\times}\) smaller} (${\sim}4$M vs.\ 110M parameters).
\LightRet{} requires no tokenizer, processes arbitrary Unicode text, and runs
efficiently on CPU-constrained deployments.
\end{abstract}
